% Post-CP26 theorem; exact independently audited candidate retained.
\subsection{Nonproduct sources of every odd polarized multiplier}
\label{subsec:nonordinary-arbitrary-odd-neighbors-candidate}
\status{Proved and independently checked.}

\begin{lemma}[The connected principal complement and its degree]
\label{lem:nonordinary-mixed-principal-complement-candidate}
\textup{[Proved; independent audits passed.]}
Let $(A,\lambda)$ be an indecomposable principally polarized abelian
surface over an algebraically closed field of characteristic two,
with two-rank one. If $A$ is not simple, it contains exactly one
ordinary elliptic subvariety $E_o$ and exactly one supersingular
elliptic subvariety $E_s$. For their inclusions $i_o,i_s$, the maps
\[
 p_o=i_o^\vee\lambda:A\longrightarrow E_o^\vee,
 \qquad p_s=i_s^\vee\lambda:A\longrightarrow E_s^\vee
\]
have respective scheme-theoretic kernels $E_s,E_o$, which are smooth
and connected. The two positive intersection degrees agree:
\[
 d=\Theta\cdot E_o=\Theta\cdot E_s>1,
 \qquad E_o\cap E_s=E_o[d]=E_s[d]
\]
as subgroup schemes inside $A$.
\end{lemma}

\begin{proof}
We first justify the connected-kernel duality used here, including
its scheme structure. For an elliptic inclusion $i:E\hookrightarrow A$,
the dual $p=i^\vee:A^\vee\to E^\vee$ is surjective: composing it
with any polarization of $A$ and $i$ gives the nonzero polarization
on $E$. Put $C=(\ker p)^0_{\mathrm{red}}$, an abelian subvariety,
and form $q:A^\vee\to D=A^\vee/C$. Then $p=u q$ for an isogeny
$u:D\to E^\vee$. The quotient has geometrically connected smooth
proper fibres, and $q_*\mathcal O=\mathcal O$ universally. Projection
formula therefore shows that a rigidified line bundle whose pullback
by $q$ is trivial is itself trivial, over every base. Thus $q^\vee$
has trivial scheme-theoretic kernel and is a closed immersion.
Dualizing $p=u q$ gives
\[
                       i=q^\vee u^\vee.
\]
Since $i$ has trivial kernel, $u^\vee$ does too. Isogeny duality gives
$\deg u=\deg u^\vee=1$, hence $u$ is an isomorphism. Consequently
$\ker p=C$, smoothly and scheme-theoretically. This supplies the
full connected-kernel assertion, rather than identifying only its
reduced identity component. The quotient and duality inputs are
\cite[Proposition~5.31, Theorem~7.5, Exercise~7.7, and
Theorem~12.2 with Remark~12.3]{EdixhovenVanDerGeerMoonenAV}; the
projection-formula argument above also proves the injectivity
required in that exercise.

If $A$ is not simple, Poincare splitting supplies two elliptic
isogeny factors. Their two-ranks sum to one, so one is ordinary and
the other supersingular. There is at most one elliptic subvariety
of either type: two distinct ones have finite intersection and
their addition map is an isogeny onto $A$; two ordinary factors
would give two-rank two, and two supersingular factors two-rank zero.
Denote the unique subvarieties by $E_o,E_s$.

There is no nonzero homomorphism between an ordinary and a
supersingular elliptic curve, since any nonzero homomorphism would
be an isogeny and preserve two-rank. Hence $p_o$ vanishes on $E_s$.
The connected-kernel result and dimension one show that its entire
kernel is $E_s$. Similarly $\ker p_s=E_o$. If
$d_i=\Theta\cdot E_i$, the restricted polarization
$i_i^\vee\lambda i_i$ is $d_i$ times the principal elliptic
polarization. Therefore
\[
 E_o\cap E_s=\ker(p_o i_o)=E_o[d_o],
 \qquad E_o\cap E_s=\ker(p_s i_s)=E_s[d_s].
\]
Taking finite group-scheme lengths gives $d_o^2=d_s^2$, hence
$d_o=d_s=d$. If $d=1$, the retraction
$r=(i_o^\vee\lambda i_o)^{-1}i_o^\vee\lambda$ splits $A$ as
$E_o\times\ker r$. The two factors are orthogonal for $\lambda$,
which becomes their product principal polarization. This contradicts
indecomposability, so $d>1$.
\end{proof}

\begin{theorem}[A nonproduct incoming source for every odd multiplier]
\label{thm:nonordinary-arbitrary-odd-neighbor-candidate}
\textup{[Proved; independent audits passed.]}
Let $(A,\Theta_A)$ be any indecomposable nonordinary principally
polarized surface over an algebraically closed field of
characteristic two. For every positive odd integer $n$, there is an
indecomposable principally polarized surface $(B,\Theta_B)$ and an
etale isogeny $\phi:B\to A$ satisfying
\[
 \phi^*\lambda_{\Theta_A}=n\lambda_{\Theta_B},
 \qquad \deg\phi=n^2,\qquad f(B)=f(A).
\]
For each odd prime $\ell$, at most one maximal isotropic subgroup
of $A[\ell]$ has a decomposable principally polarized quotient;
in two-rank zero none does.
\end{theorem}

\begin{proof}
Fix an odd prime $\ell$. For a maximal isotropic subgroup
$H\subset A[\ell]$ the principal polarized quotient exists and has
\[
 q:A\longrightarrow B=A/H,\qquad
 q^*\lambda_{\Theta_B}=\ell\lambda_{\Theta_A},\qquad
 |H|=\ell^2.
\]
This is exactly the polarized quotient construction already used
in Lemma~\ref{lem:root-nonordinary-nonproduct-three-neighbor-candidate},
with three replaced by $\ell$; the perfect alternating Weil pairing
is on $A[\ell]\simeq\mathbf F_\ell^4$. Both $q$ and its polarized
dual are etale and induce isomorphisms on the two-divisible groups.

If $f(A)=0$, Cartier on invariant one-forms has rank one: the
smooth genus-two theta normalization has forms $dx,x\,dx$ with
Cartier images $0,dx$. A decomposable two-rank-zero principal
surface is a product of supersingular elliptic curves and has
Cartier rank zero. The etale pullback of invariant differentials
is an isomorphism commuting with Cartier, so $B$ cannot be such
a product. This is the same rank argument as for the proved
three-neighbour theorem and is independent of $\ell$.

Suppose $f(A)=1$ and $B$ is decomposable, with ordinary and
supersingular elliptic axes $E_o',E_s'$. Their connected inverse
images in $A$ are elliptic subvarieties of the corresponding types.
Thus $A$ is not simple, and they are exactly the unique $E_o,E_s$
of the preceding lemma. The restrictions $q|_{E_i}$ are isogenies
onto $E_i'$. Restricting the pullback polarization yields
\[
                    \deg(q|_{E_i})=\ell d.
\]
Their kernels are subgroup schemes of $E_i[\ell]$, so these degrees
divide $\ell^2$. Since $d>1$, necessarily $d=\ell$ and
$\deg(q|_{E_i})=\ell^2$. Consequently
\[
                         H=E_o[\ell]=E_s[\ell].
\]
The possible bad subgroup is therefore unique. A symplectic basis
of $A[\ell]$ exhibits at least two maximal isotropic subgroups,
namely the spans of its two $e$ vectors and its two $f$ vectors.
Choose a subgroup different from the possible bad one. Its
principal quotient is indecomposable.

Now factor $n=\ell_1\cdots\ell_s$ with repetitions. Successively
choose such nonproduct outgoing principal neighbours, retaining
two-rank at each step. Their composition $q:A\to B$ has polarized
multiplier $n$ and degree $n^2$. Its polarized dual
\[
           \phi=\lambda_{\Theta_A}^{-1}q^\vee\lambda_{\Theta_B}
                     :B\longrightarrow A
\]
satisfies $\phi q=[n]_A$ and $q\phi=[n]_B$, and hence the stated
pullback formula. Its kernel lies in $B[n]$, so it is etale.
For $n=1$, take the identity. All constructions use finite algebraic
data and thus descend to a finite extension of any field of
definition; no isogeny over that original field is asserted.
\end{proof}

\begin{corollary}[All odd multiples of three on fixed nonordinary surfaces]
\label{cor:nonordinary-all-odd-multiples-three-candidate}
\textup{[Proved; independent audits passed.]}
Let $(A,M)$ be a principally polarized surface over a perfect field
of characteristic two, with $f=f(A)\in\{0,1\}$, and put $L=M^3$.
For every positive odd $d$ divisible by three, the geometric and
arithmetic relative groups of $|L^d|$ are
\[
                         \mathrm{GL}_{9d^2-1}(\mathbf F_2).
\]
For every positive even $d$ with $d\equiv0$ or $4\pmod6$, and also
for $d=2$, both groups are
\[
 \operatorname{Hom}(\mathbf F_2^{1+f},\mathbf F_2^{9d^2-4})
                  \rtimes\mathrm{GL}_{9d^2-4}(\mathbf F_2).
\]
In two-rank zero every even degree was already proved by
Corollary~\ref{cor:f0-all-even-nine-theta-transfer-candidate}.
\end{corollary}

\begin{proof}
For a nonproduct geometric target, take $n=d/3$ in the incoming
neighbour theorem and apply
Theorem~\ref{thm:nine-theta-arbitrary-odd-norm-candidate}; for $n=1$
use the previously proved fixed degree-three theorem. The source
is nonproduct and nonordinary, exactly as that norm theorem requires.
Geometrically decomposable nonordinary surfaces are the already
classified supersingular or mixed product cases.

If $d\equiv0\pmod6$, write $d=3+(d-3)$; both summands are odd
multiples of three. If $d\equiv4\pmod6$, write $d=1+(d-1)$,
using the proved degree-one theorem. The even-addition theorem
then gives the displayed full pointwise stabilizer of the
$1+f$ invariant covectors. Degree two is its separate proved base.
In odd degree geometric fullness already forces arithmetic
fullness. In even degree the entire geometric invariant-covector space
descends pointwise over the perfect base field, by
Theorem~\ref{thm:nonordinary-ppav-all-covectors}, giving the
same equality. Translation
identifies any representing principal bundle with the geometric
symmetric representative, so the statement applies to the stated
$M$ and its actual complete systems.
\end{proof}

\begin{remark}[Residual fixed-surface gates]
\label{rem:nonordinary-arbitrary-odd-neighbor-scope-candidate}
\textup{[Proved; independent audits passed.]}
The quotient argument concerns existence of a nonproduct source,
not nodality of its individual theta image. Its use is justified
by the separate proved nine-theta transfer. In two-rank one it
does not settle degree five or degree eight, or in general the
remaining even class $2\pmod6$ beyond the degree-two base. In
two-rank zero degree eleven is settled by
Corollary~\ref{cor:f0-odd-multiples-eleven-candidate}; the first
remaining odd degree is thirteen.
The proof is symbolic and invokes no new finite census. The complete complement/duality, quotient, and composition
arguments passed the parent full proof/source audit at
22:41:22 UTC and the obstruction/crystal audit at 22:39:59 UTC
on 19 November 2026; the latter also checked the final
pointwise arithmetic-descent precision. No numerical run
was needed or claimed for this symbolic proof.
\end{remark}
